\documentclass[11p]{article}
% Packages
\usepackage{amsmath}
\usepackage{graphicx}
\usepackage[swedish]{babel}
\usepackage[
    backend=biber,
    style=authoryear-ibid,
    sorting=ynt
]{biblatex}
\usepackage[utf8]{inputenc}
\usepackage[T1]{fontenc}
\usepackage{blindtext}
\usepackage{textcomp}
%Källor
\addbibresource{mall.bib}
\graphicspath{ {./images/} }

\title{Labrapprt \\ \small Fysik 1}
\author{Magnus Silverdal }
\date{\today}

\begin{document}

    \begin{titlepage}
        \begin{center}
            \vspace*{1cm}

            \Huge
            \textbf{Laboration 5}

            \vspace{0.5cm}
            \LARGE
            Ellära

            \vspace{1.5cm}

            \textbf{Alvin Högdal}

            \vfill


            Fysik 1

            \vspace{0.8cm}

            \includegraphics[width=0.4\textwidth]{../images/NTI Gymnasiet_Symbol_print_svart.png}

            \Large
            Teknikprogrammet\\
            NTI Gymnasiet\\
            Umeå\\
            \today

        \end{center}
    \end{titlepage}
    \section{Syfte och frågeställning}
    \paragraph{Beräkna hur stor resistansen är på ett okänt motstånd.}
    
    \section{Del 1}
    \subsection{Material och metod}
       \paragraph{Material:}
       \begin{itemize}
           \item Ett batteri
            \item En multimeter
            \item Två kablar
            \item En kopplingsplint
            \item Två krokodilklämmor
            \item Ett motstånd
       \end{itemize}
        \paragraph{Metod:}
        Parallellkoppla voltmätaren och ett motstånd.
        Antäckna resultatet.
        Serikoppla montsåndet och en ampärmätare.
        Antäckna resultatet.
        \begin{figure}[!h]
        \includegraphics[width=0.8\textwidth]{../images/img2.png}
        \caption{Exempel på en paralellkopling.}
        \label{fig:paralell}
        \end{figure}
        \paragraph{}
        \begin{figure}[!h]
        \includegraphics[width=0.5\textwidth]{../images/img1.png}
        \caption{Exempel på en seriekoppling.}
        \label{fig:serie}
        \end{figure}
    \subsection{Resultat}
    Voltmätaren gav ett resultat på 8,45V och ampärmätaren gav ett resultat på 0,08A\@.
    Vi kan nu beräkna resistansen enligt Ohms lag.
    \begin{equation}
        R = \frac{U}{I}
    \end{equation}
    Detta ger resultatet $\sim$ 100\Omega.

    \subsection{Analys}
    Resultatet 100\Omega\text{ stämmer eftersom om man kollar upp vad som är markerat på resistorn så står det 100}\Omega.
    \newpage
    \section{Del 2}
    \subsection{Material och metod}
    \paragraph{Material:}
    \begin{itemize}
        \item Ett batteri
        \item En multimeter
        \item Fyra kablar
        \item Tre kopplingsplintar
        \item Två krokodilklämmor
        \item Två motstånd
    \end{itemize}
    \paragraph{Metod:}
    Seriekoppla två motstånd.
    Använd multimetern för att bestämma strömen och spänningen över ett av motstånden.
    Antäckna resultaten.

    \subsection{Resultat}
    Multimetern gav ett resultat på 4,2V och 0,04A.
    \begin{equation}
        R = \frac{4,2}{0,04}
    \end{equation}
    Detta ger resultatet $\sim$ 100\Omega.
    \subsection{Analys}
    Båda motstånden är samma och om man kollar upp vad som står på dom så är det 100\Omega.\text{ Detta stämmer med resultated av 100}\Omega


\end{document}
